%%%%%%%%%%%%%%%%%%%%%%%%%%%%%%%%%%%%%%%%%%%%%%%%%%%%%%%%%%%%%%%%%%%%%%%%%%%
\chapter{RTEB and Generalization}
\label{ch:rteb}
%%%%%%%%%%%%%%%%%%%%%%%%%%%%%%%%%%%%%%%%%%%%%%%%%%%%%%%%%%%%%%%%%%%%%%%%%%%

If retrieval is your main use case, RTEB is more relevant than the overall MTEB aggregate.
RTEB stands for \textbf{Retrieval Text Embedding Benchmark}. The Hugging Face launch post was
published on October~1, 2025, by the MTEB team.

%%%%%%%%%%%%%%%%%%%%%%%%%%%%%%%%%%%%%%%%%%%%%%%%%%%%%%%%%%%%%%%%%%%%%%%%%%%
\section{What RTEB Is Trying to Fix}
%%%%%%%%%%%%%%%%%%%%%%%%%%%%%%%%%%%%%%%%%%%%%%%%%%%%%%%%%%%%%%%%%%%%%%%%%%%

RTEB was introduced because broad public retrieval leaderboards have a known problem: models
get repeatedly tested against the same public datasets, and training data can overlap with
benchmark data. Leaderboard gains can partly reflect benchmark familiarity rather than true
retrieval generalization.

The RTEB authors frame this as a generalization gap between public benchmark performance and
performance on new, unseen retrieval problems. The benchmark is explicitly trying to be:

\begin{itemize}
    \item Retrieval-first
    \item More realistic
    \item More resistant to benchmark overfitting
\end{itemize}

The Hugging Face post makes a useful distinction between a \textit{leaderboard score} and a
\textit{zero-shot score}. A model can look strong on a familiar public benchmark while still
generalizing poorly to new retrieval tasks.

%%%%%%%%%%%%%%%%%%%%%%%%%%%%%%%%%%%%%%%%%%%%%%%%%%%%%%%%%%%%%%%%%%%%%%%%%%%
\section{Why RTEB Matters More Than Generic MTEB for Local RAG}
%%%%%%%%%%%%%%%%%%%%%%%%%%%%%%%%%%%%%%%%%%%%%%%%%%%%%%%%%%%%%%%%%%%%%%%%%%%

Your actual local use case is not broad embedding capability across classification, clustering,
STS, and bitext mining. It is: can this model retrieve the right chunks for my queries?

That makes RTEB a better decision aid because it focuses on retrieval quality, domain realism,
and multilingual and specialized search settings.

%%%%%%%%%%%%%%%%%%%%%%%%%%%%%%%%%%%%%%%%%%%%%%%%%%%%%%%%%%%%%%%%%%%%%%%%%%%
\section{The Core Design: Open + Private Datasets}
%%%%%%%%%%%%%%%%%%%%%%%%%%%%%%%%%%%%%%%%%%%%%%%%%%%%%%%%%%%%%%%%%%%%%%%%%%%

RTEB uses a hybrid benchmark design:

\begin{itemize}
    \item \textbf{Open datasets} --- fully public and reproducible
    \item \textbf{Private datasets} --- evaluated by the maintainers to test generalization
          to unseen data
\end{itemize}

This is the most important thing to understand about RTEB. Strong public-only performance can
be misleading. If a model drops sharply on the private side, that suggests benchmark
overfitting or weak generalization.

For transparency, the maintainers provide descriptive statistics, dataset descriptions, and
sample (query, document, relevance) triplets. The benchmark balances reproducibility with
anti-overfitting protection.

%%%%%%%%%%%%%%%%%%%%%%%%%%%%%%%%%%%%%%%%%%%%%%%%%%%%%%%%%%%%%%%%%%%%%%%%%%%
\section{Default Metric: nDCG@10}
%%%%%%%%%%%%%%%%%%%%%%%%%%%%%%%%%%%%%%%%%%%%%%%%%%%%%%%%%%%%%%%%%%%%%%%%%%%

The RTEB material explicitly states that the default leaderboard metric is \texttt{nDCG@10}.
That is a strong choice for retrieval because it emphasizes whether relevant documents appear
near the top---which is exactly what matters when a RAG pipeline passes only the top 5 or
top 10 chunks into the final prompt.

%%%%%%%%%%%%%%%%%%%%%%%%%%%%%%%%%%%%%%%%%%%%%%%%%%%%%%%%%%%%%%%%%%%%%%%%%%%
\section{What Kinds of Retrieval RTEB Covers}
%%%%%%%%%%%%%%%%%%%%%%%%%%%%%%%%%%%%%%%%%%%%%%%%%%%%%%%%%%%%%%%%%%%%%%%%%%%

RTEB is not one dataset. It is a retrieval benchmark family spanning multiple domains and
languages. The benchmark currently covers 20 languages, with datasets organized into groups
by domain and language. One dataset can belong to multiple groups---a German legal dataset
counts as both \texttt{german} and \texttt{legal}.

Domain coverage includes:

\begin{itemize}
    \item \textbf{Legal:} AILACasedocs, AILAStatutes, LegalSummarization, LegalQuAD
    \item \textbf{Finance:} FinanceBench, HC3Finance, FinQA
    \item \textbf{Code:} HumanEval, MBPP, APPS, DS1000, WikiSQL
    \item \textbf{Healthcare:} ChatDoctor\_HealthCareMagic, HC3 Medicine, Cure, TripClick
    \item \textbf{Multilingual and general:} MIRACLHardNegatives, JaQuAD, FreshStack
\end{itemize}

RTEB benchmark names in the catalog include:

\begin{itemize}
    \item \texttt{HumanEvalRetrieval}, \texttt{MBPPRetrieval}, \texttt{WikiSQLRetrieval}
    \item \texttt{FreshStackRetrieval}, \texttt{FinanceBenchRetrieval}
    \item \texttt{ChatDoctorRetrieval}, \texttt{MIRACLRetrievalHardNegatives}
\end{itemize}

This is important because retrieval quality is highly domain-sensitive. A model that looks
good on generic QA-style retrieval may not be the best model for code search, technical docs,
multilingual knowledge bases, or legal corpora.

%%%%%%%%%%%%%%%%%%%%%%%%%%%%%%%%%%%%%%%%%%%%%%%%%%%%%%%%%%%%%%%%%%%%%%%%%%%
\section{How to Read RTEB for Your Use Case}
%%%%%%%%%%%%%%%%%%%%%%%%%%%%%%%%%%%%%%%%%%%%%%%%%%%%%%%%%%%%%%%%%%%%%%%%%%%

You should not ask: ``which model is \#1 on RTEB overall?''

You should ask: ``which model performs well on the RTEB slice closest to my workload?''

\begin{itemize}
    \item If your corpus is mostly English docs, \texttt{RTEB(eng, beta)} matters more than
          the multilingual aggregate
    \item If your main problem is code retrieval, the code-heavy tasks matter more than
          finance or legal tasks
    \item If your content spans multiple languages, the multilingual and language-specific
          slices matter more than English-only results
\end{itemize}

\textbf{RTEB slices relevant to local infrastructure workloads:}

\begin{description}
    \item[\texttt{RTEB(beta)}] Broad retrieval signal across domains
    \item[\texttt{RTEB(eng, beta)}] Best general-purpose check if most material is English
    \item[Code-heavy tasks] For source code, documentation, or engineering notes
    \item[Multilingual retrieval tasks] If knowledge base or queries are not strictly English
\end{description}

Language-specific subsets include \texttt{RTEB(deu, beta)}, \texttt{RTEB(fra, beta)}, and
\texttt{RTEB(jpn, beta)}.

%%%%%%%%%%%%%%%%%%%%%%%%%%%%%%%%%%%%%%%%%%%%%%%%%%%%%%%%%%%%%%%%%%%%%%%%%%%
\section{What RTEB Still Does Not Solve}
%%%%%%%%%%%%%%%%%%%%%%%%%%%%%%%%%%%%%%%%%%%%%%%%%%%%%%%%%%%%%%%%%%%%%%%%%%%

RTEB is better aligned with retrieval than the broad MTEB aggregate, but it is not perfect.
The official launch notes several current limitations:

\begin{itemize}
    \item It is still in beta
    \item It is text-only
    \item Language coverage is still expanding
    \item About half of the current retrieval datasets are repurposed from QA datasets
\end{itemize}

QA-derived retrieval can favor lexical overlap more than true production retrieval sometimes
does. So even with RTEB, your own internal retrieval eval still matters.

%%%%%%%%%%%%%%%%%%%%%%%%%%%%%%%%%%%%%%%%%%%%%%%%%%%%%%%%%%%%%%%%%%%%%%%%%%%
\section{What RTEB Shows About Qwen3-Embedding}
%%%%%%%%%%%%%%%%%%%%%%%%%%%%%%%%%%%%%%%%%%%%%%%%%%%%%%%%%%%%%%%%%%%%%%%%%%%

Qwen3-Embedding does not have published first-party RTEB scores as of this writing. That
absence is the core open question in the original model recommendation: MTEB rank-1 without
a corresponding RTEB result means the generalization question is unanswered, not confirmed.

Indirect evidence:

\begin{itemize}
    \item Octen-Embedding-8B, fine-tuned on Qwen3-8B base, scores \textbf{0.8045 on RTEB
          overall} and performs slightly higher on private data than public (0.8157 private
          vs 0.7953 public). That inverted gap is a strong generalization signal and suggests
          the Qwen3 foundation model is not contaminated.
    \item Qwen3-Embedding's three-stage training (contrastive pretraining on weakly supervised
          data $\to$ supervised fine-tuning $\to$ model merging) is designed for domain breadth,
          which is consistent with generalization.
    \item Qwen3-Embedding outperforms BGE-M3 by \textasciitilde18\% on MTEB English v2 and
          \textasciitilde8\% on MMTEB, a larger margin than benchmark contamination alone
          would typically produce.
\end{itemize}

Tentative interpretation: Qwen3-Embedding is probably not severely contaminated, but
``probably not'' is not the same as confirmed.

\textbf{RTEB-confirmed alternatives:}

\begin{center}
\begin{tabular}{lll}
\toprule
\textbf{Model} & \textbf{RTEB score} & \textbf{Notes} \\
\midrule
Octen-Embedding-8B & 0.8045 & Fine-tuned on Qwen3-8B; best published RTEB result \\
Voyage-3-large & 0.7812 & API-only; no local serving \\
Gemini-embedding-001 & 0.7602 & API-only \\
BGE-M3 & not published & Strong multilingual; GGUF available \\
\bottomrule
\end{tabular}
\end{center}

For local llama.cpp serving specifically, BGE-M3 is the strongest RTEB-adjacent alternative:
it supports dense, sparse (SPLADE), and ColBERT-style retrieval in a single model, has a
well-established generalization reputation predating the MTEB contamination era, and 8{,}192
token context.

%%%%%%%%%%%%%%%%%%%%%%%%%%%%%%%%%%%%%%%%%%%%%%%%%%%%%%%%%%%%%%%%%%%%%%%%%%%
\section{Benchmark-Reading Order for Local RAG}
%%%%%%%%%%%%%%%%%%%%%%%%%%%%%%%%%%%%%%%%%%%%%%%%%%%%%%%%%%%%%%%%%%%%%%%%%%%

\begin{enumerate}
    \item Domain-relevant RTEB retrieval results
    \item Broader retrieval metrics like MTEB retrieval subsets
    \item Overall MTEB aggregate
    \item Secondary task families like clustering or classification (background signal only)
\end{enumerate}

\begin{notebox}[Before finalising a model choice]
Check whether Qwen3-Embedding RTEB scores have been published since this was written (the
RTEB leaderboard is live and updated). Run a small domain-specific internal eval against your
actual retrieval corpus --- even 50--100 query-document pairs with manual relevance labels
will reveal gross distribution mismatch.
\end{notebox}
